% `tex.web` §1153 and §1136: a `{` in math mode is a *subformula*, and it gets
% a nest level of its own.
%
%     mmode+left_brace: begin tail_append(new_noad); back_input;
%       scan_math(nucleus(tail));
%       end;
%
%     procedure push_math(@!c:group_code);
%     begin push_nest; mode:=-mmode; incompleat_noad:=null;
%     new_save_level(c);
%     end;
%
% `mode := -mmode` — inner — whatever style the material will eventually be
% set in. Only the outermost level of a *display* is `mmode` positive; `$…$`
% and every braced group inside either are not. So `$$B{C}D$$` reports
% `{math mode: the letter C}` and then `{display math mode: the letter D}`.
%
% We read the *style* to decide, and the style is `display_style` all the way
% down a display — so every subformula in one was named for the display it sat
% in. §211's mode name is not decoration: it is how a reader of a
% `\tracingcommands` log tells which list a command went to, and §218's
% `\showlists` heading is built from the same answer.
%
% The style and the mode are genuinely different questions, which is why one
% cannot stand in for the other: a `{}` group in a display is inner mode and
% display style, and `\hbox{$Q$}` inside a display is *text* style and inner
% mode both.
\catcode`\{=1 \catcode`\}=2 \catcode`\$=3 \catcode`\^=7 \catcode`\_=8
\tracingonline=1 \tracingcommands=2
\hsize=100pt \parindent=0pt
\font\rm=cmr10 \textfont0=\rm \scriptfont0=\rm \scriptscriptfont0=\rm \rm
\font\mi=cmmi10 \textfont1=\mi \scriptfont1=\mi \scriptscriptfont1=\mi
\font\sy=cmsy10 \textfont2=\sy \scriptfont2=\sy \scriptscriptfont2=\sy
\font\ex=cmex10 \textfont3=\ex \scriptfont3=\ex \scriptscriptfont3=\ex
\showthe\escapechar

% Nested groups: in and out, one level at a time.
\setbox0\vbox{\noindent A$$B{C{D}E}F$$G\par}

% Scripts open one the same way — §1176's `scan_math` on the field.
\setbox1\vbox{\noindent A$$B^{C}_{D}E$$F\par}

% ...and so do the two halves of a fraction, when §1184 splits the list.
\setbox2\vbox{\noindent A$$B\over C$$D\par}

% \vcenter is a vertical list inside a formula, and what follows it is back in
% the display.
\setbox3\vbox{\noindent A$$\vcenter{\hbox{x}}B$$C\par}

% An \hbox in a display holds a formula of its own: inner mode, text style.
\setbox4\vbox{\noindent A$$\hbox{$Q$}B$$C\par}

% Inline math is inner already, so a group in one changes nothing.
\setbox5\hbox{A$B{C}D$E}
\showthe\escapechar
\end
